\chapter{Conclusiones y trabajo futuro}\label{cap:conclusiones}

\section{Conclusiones}

El trabajo se propuso estudiar tres dimensiones, y hay una conclusión por cada una.

\paragraph{Armonía: no encaja en ningún estilo.} La banda toma elementos de varios estilos, pero no encaja en ninguno. El
37.5\,\% de los casos queda sin vecino, frente al 5.2\,\% de las piezas reales de rock, y ninguna canción pasa de 17.7
cuando una pieza típica daría 50. Las fundamentales son de rock modal, sin dominante, pero desplazadas hacia \fl VII y
\fl III, a una zona que ningún corpus ocupa.

\paragraph{Melodía: lo nuevo.} Dos grados en cuarto de tono concentran entre el 48 y el 82\,\% de los cuartos de tono de
cada canción. Son notas fijas de la escala, no adornos, y forman dos familias. La frigia (\fl2+, \fl6+) es la más común en
el radif persa y en King Gizzard. La dórica (7+, 6+, \fl3+) no aparece en ningún corpus medido.

\paragraph{Síntesis: un híbrido.} Hay ciclos de 5 y de 7 tiempos, como dice la banda; la armonía se parece a la del rock
modal; los cuartos de tono se reparten como en el radif; y el vecino más cercano es King Gizzard.

\medskip
La respuesta a la pregunta del trabajo es que la música de Angine de Poitrine no corresponde a un estilo existente. Es una
combinación nueva de lenguajes que ya existían, con un elemento propio: los cuartos de tono como notas fijas de la escala.

Además del resultado, el trabajo deja una herramienta abierta que transcribe y analiza en 24-EDO desde el audio, y un
método de comparación con corpus que aplica a los corpus el error de la propia herramienta y está calibrado para no
exagerar la novedad.

\section{Trabajo futuro}

El siguiente paso es un tabroll que mejore solo, para cualquier canción, mediante un ciclo de cuatro pasos:
\begin{enumerate}
  \item sintetizar partituras de Guitar Pro con cuartos de tono;
  \item entrenar con ellas un transcriptor en 24-EDO;
  \item transcribir canciones nuevas;
  \item validar con referencias de oído.
\end{enumerate}
Así habría datos sin límite, sin transcribir a mano. Además:
\begin{itemize}
  \item \textbf{Melodía:} tratar las vueltas del looper como varias tomas de la misma frase, para separar las capas.
  \item \textbf{Armonía:} elegir el centro tonal de forma robusta en cualquier canción.
  \item \textbf{Datos:} un corpus de maqam árabe y el acceso a ProgGP, para comprobar el compás cambiante frente al rock
  progresivo.
  \item \textbf{Producto:} un tabroll web interactivo, el análisis de otros músicos microtonales, como King Gizzard y
  Sevish, y su uso en la enseñanza.
\end{itemize}
